% 本文件由 LaTeX 排版 Agent 根据有效 translation.json 自动生成。
% author=Tertullian; work_id=0313; title=驳赫莫杰尼斯

\chapter{驳赫莫杰尼斯}

% structure: p0001 source=h1 target=omitted reason=page-title-already-rendered-by-parent

% structure: p0002 source=h2 target=section reason=relative-heading-rank; base=h2

\section{第一章 赫莫杰尼斯的意见：藉古时的规定法则显明为异端；并非源自基督教，而是来自异教哲学；并提及若干主张。}

我们惯于为了缩短论证，针对异端者的\Emph{晚近}日期设定规则。因为按照我们的规则，真理享有优先权，真理也曾预言会有异端；所以，所有后来的意见都必须预先被判定为异端，因为它们正是被更古老的真理规则预言将要发生的。而现在，赫莫杰尼斯的学说就带有这种新颖性的污点。简言之，他是当今活在世上的人；他本性就是个异端者，且暴躁不安，误以为饶舌是口才，把厚颜当作坚定，并认为诋毁个人是良心责任。此外，他在绘画中藐视天主的法律，坚持多次婚姻，以天主的法律为情欲辩护，却又在艺术方面藐视它。他以双重方式作假——用他的烙笔和笔。他在教义上和肉体上都是彻底的奸夫，因为他确实沾染了你们那些婚姻贩子的污秽，并且也未能持守信德规则，正如宗徒本人的赫莫杰尼斯那样。\BibleRef{弟后 1:15} 然而，不必在意这人，我质疑的是他的学说。他似乎不承认别的基督为主，尽管他以不同方式持有祂；但正是这种信仰上的差异使他实际上将祂变成另一位——不，他夺走了祂一切属于天主的属性，因为他不愿承认祂从无中创造了万物。因为，他离开基督徒转向哲学家，从教会转向\OriginalTerm{学园}{Academy}和\OriginalTerm{画廊}{Porch}，在那里从斯多亚派学得将\OriginalTerm{物质}{Matter}与主置于同等地位，仿佛它也是非受生、非受造的，根本没有起始和终结；据他说，主以后从物质中创造了万物。\par

% structure: p0004 source=h2 target=section reason=relative-heading-rank; base=h2

\section{第二章 赫莫杰尼斯在从纯粹的异端假设进行乖谬归纳之后，结论说天主从已存的物质中创造了万物。}

我们这位很糟糕的画家，用这样的论证涂上了他这第一道完全无光的色调：他首先设定前提，说主造万物，或是出于自己，或是出于无，或是出于某物；以便在证明了祂不可能出于自己或出于无造物之后，由此肯定剩下的命题，即祂出于某物造物，因此那某物就是物质。他说，祂不能出于自己造万物；因为主出于自己所造的任何事物都会是祂的部分；但祂不可分割为部分，因为作为主，祂是不可分、不变、永远相同的。此外，如果祂曾出于自己造了什么，那就会是祂的某物。然而，凡被造的以及祂所造的，都必被视为不完全，因为它是出于一部分而造的，而且祂是出于一部分造了它；或者，如果祂所造的是整体（而祂自己是整体），那么祂就必须同时是整体又不是整体；因为祂必须是整体才能产生自己，又不是整体才能从自己产生。但这是最困难的立场。因为如果祂存在，祂就不能被造，因为祂已存在；如果祂不存在，祂就不能造，因为祂是非存在。他\Emph{坚持}说，那永恒存在者并不\Emph{进入}存在，而是永远存在。于是他结论说，祂没有出于自己造任何东西，因为祂从未进入那种可能使自己出于自己造物的状态。同样，他主张祂也不能出于无造万物——如此：他把主定义为善的，甚至至善的，祂必须意愿造出与祂自己一样善而卓越的事物；的确，祂不可能意愿或造出任何不善的、甚至不是至善的事物。因此，万物应当被祂造得善而卓越，符合祂自己的状况。然而，经验表明，甚至恶的事物也是祂造的；当然，不是出于祂自己的意愿和喜悦；因为如果是出于祂自己的意愿和喜悦，祂一定不会造任何不合宜或不配祂自己的事物。因此，祂不是出于自己意愿所造的，必须理解为是由于某物的缺陷而造的，而那就是\OriginalTerm{物质}{Matter}，毫无疑问。\par

% structure: p0006 source=h2 target=section reason=relative-heading-rank; base=h2

\section{第三章 赫莫杰尼斯的一个论证。答复：虽然天主是永远适用于神性存在的称号，主和父只是相对称谓，并非永远适用。指出赫莫杰尼斯论证中的矛盾。}

他又提出了另一点：既然天主永远是天主，就绝没有一个时期天主不是主。但是，祂不可能像永远是天主那样，永远被视为主，除非在过去的永恒中，始终存在某种东西，祂可以被视为其永远的主。于是他得出结论说，天主始终有物质与祂共存，作为其主。我将立即撕开他编织的这一套。我之所以愿意将其如此陈明，是为了使不熟悉这一问题的人了解，从而知道他的其他论点也只需被理解，就能被驳斥。那么，我们肯定说，\Emph{天主}这个名称始终与祂自身并在祂自身内存在——但\Emph{主}并非如此永恒。因为一方的情况与另一方不同。天主是实体本身，即天主性的名称；但主不是实体的名称，而是权能的名号。我\Emph{主张}实体始终以其自己的名称存在，即天主；\Emph{主的称号}是后来加上去的，作为某种附加之事的标志。因为从那些要受主权作用的受造物开始存在的那一刻起，\Emph{天主}由于加入那权能，既成了主，也接受了它的名号。因为天主也是父，也是审判者；但祂并不总是父和审判者，仅因为祂总是天主。因为祂不可能在圣子之前成为父，也不可能在罪之前成为审判者。然而，曾有一段时间，罪和圣子都不与祂同在；罪将构成使主成为审判者，圣子则构成父。这样，在祂要成为其主的事物存在之前，祂并不是主。祂只是在将来的某个时刻才成为主：正如祂因圣子成为父，因罪成为审判者，同样，祂也因祂所造、为要服事祂的事物而成为主。\Person{Hermogenes}，你是否觉得我在编织论点？圣经何等巧妙地帮助我们，它有区别地将这两个称号应用于祂，并在适当的时候揭示它们！因为\Emph{天主}这个永远属于祂的称号，圣经在最初就提到：\BibleQuote{在起初天主创造了天地；}\BibleRef{创 1:1}当祂继续逐一创造那些祂将成为主的受造物时，圣经仅仅提到天主。\BibleQuote{天主说}，\BibleQuote{天主造了}，\BibleQuote{天主看了}；但我们至今无处找到\Emph{主}。但当祂完成了全部创造，特别是造了人——人要以一种特殊的方式理解祂的主权——那时祂才被称为主。于是\Emph{圣经}加上了\Emph{主}的名称：\BibleQuote{然后上主天主，\Emph{Deus Dominus}，取了祂所形成的人；}\BibleRef{创 2:15}\BibleQuote{上主天主命令了亚当。}\BibleRef{创 2:16}从此，那先前只是天主的，成为\Emph{主}，当祂有了可以成为其主的事物时。因为对祂自己，祂永远是天主；但对万物，祂只是在成为主的时候才是天主。因此，既然\Person{Hermogenes}将以为物质是永恒的，理由是主是永恒的，那么同样将显明，没有任何东西存在，因为显然\Emph{作为主}并不总是存在。现在，我也要亲自为无知之人添上一句评论——\Person{Hermogenes}是其中极端的一例——并且实际上要以他自己的论点反驳他。因为当他否认物质是受生或受造时，我发现，即使在这些条件下，\Emph{主}这个称号对天主而言在涉及物质时也不合适，因为物质必定是自由的，由于没有起源就没有作者。它过去的存在不归功于任何人，因此它不能属于任何人。因此，自从天主通过从物质中创造万物而对它行使权能以来，虽然物质一直体验到天主是它的主，但毕竟，物质表明天主并非以主的关系存在，尽管祂一直是如此。\par

% structure: p0008 source=h2 target=section reason=relative-heading-rank; base=h2

\section{第四章 埃尔摩格内斯赋予物质神圣属性，从而制造了两个神}

至此，我将开始讨论\OriginalTerm{物质}{Matter}，即（按赫尔谟根尼的说法）天主如何将它与自己并列，视为同属非受生、非受造、同样永恒，且被宣示为无始无终。因为对天主的认识，除了永恒，还能有什么？永恒的状态，除了始终存在，并且因其无始无终的特权而永远存在，还能有什么？既然这是天主的属性，它就只属于天主——当然，理由是：如果它可以归给任何其他存有，那就不再是天主的属性，而是与天主一起，也属于那被归给的存有。因为\Quote{虽有称为神的}，以名称而言，\Quote{或在天上，或在地上，然而我们只有一个天主，就是圣父，万物都出于祂}；\BibleRef{格前 8:5}因此，在我们看来，那属于天主的属性理应只与天主相关，而且（正如我已说过的）当这属性被另一存有分享时，它便不再是这样的属性。既然祂是天主，那么这品质的一个独特标记必然是它只限于一位。否则，若那没有相等之物的东西不是唯一的，还有什么会是唯一和独特的？若那超越万有、先于万有、万有由之而出者不是至高的，还有什么会是至高的？因拥有这些，祂是唯一的天主；因独自拥有，祂是唯一的一位。若另一者也共有了这些，那么有多少拥有这些天主属性者，就有多少神。因此，赫尔谟根尼引入了两位神：他将物质引入作为天主的平等者。但天主必须是一位，因为至高的才是天主；除了那唯一的，没有什么是至高的；而若有任何与它相等之物，它就不可能是唯一的；物质若被视为永恒，便与天主相等。\par

% structure: p0010 source=h2 target=section reason=relative-heading-rank; base=h2

\section{第五章　赫尔谟根尼对自己的论点闪烁其词，仿佛有所畏惧。在赋予物质神圣品质之后，又试图使之低于天主}

但天主是天主，物质是物质。仿佛名称上的差异就能阻止它们平等，而它们却被声称具有相同的状态！姑且承认它们的本性不同，也假定它们的形式并不相同——只要它们绝对的状态只有一种模式，那又有什么关系？天主是非受生的；物质不也是非受生的吗？天主永远存在；物质不也是永远存在的吗？两者都无始；两者都无终；两者都是宇宙的创造者——创造它的祂，以及祂用以创造它的物质。因为既然宇宙由物质构成，就不能认为物质不是万物的创造者。他会作何回答？他会说，一旦物质有了属于天主的东西，它便不能与天主比较；因为不拥有全部（神性），它就无法完全对应比较吗？但他为天主保留了什么呢，以致他不应承认已将全部神性赋予了物质？他回答说，即使这是物质的特权，天主的权威和本质仍必须保持不变，凭此祂被视为万有的唯一和首要创造者及主宰。然而，真理坚持天主的唯一性，其方式便是强调凡属于天主自身的，都只属于祂自己。因为只有它单独属于祂，它才属于祂自身；因此，若不允许任何其他存有拥有天主的任何东西，就不可能接纳另一位神。那么，你说，按此说法，我们自己也没有任何属于天主的东西。但我们的确拥有，并将继续拥有——不过那是从祂领受，而非出于自身。因为我们甚至要成为神，若我们配得成为祂所宣称的那群\Quote{我曾说：你们是神}，以及\Quote{神站在众神的会中}的人。但这出于祂的恩宠，并非由于我们有任何属性，因为只有祂能使人成为神。然而，赫尔谟根尼却认为物质的属性就是它与天主共有的东西。否则，如果它从天主领受了属于天主的属性——我指它的永恒性——那么人们甚至可以设想它既拥有与天主共同的属性，同时又非天主。但他既承认一种属性与天主共有，又希望他不拒绝赋予物质的东西终究是天主独有特权，这是何等的矛盾！\par

% structure: p0012 source=h2 target=section reason=relative-heading-rank; base=h2

\section{第六章　赫尔谟根尼陷入困境：他神化物质，却又不愿使之与神圣创造者平等}

他宣称天主的属性依然安全地属于祂：祂是唯一的天主、是元始、是万物的创始者、万有的主，且无可比较——这些品质他立刻也归给了\OriginalTerm{物质}{Matter}。祂确实是天主。天主也必将同样作证；但祂也曾有时指着自己起誓，说没有其他神像祂一样。\BibleRef{依 45:23}然而，赫尔摩格内斯要使祂成为说谎者。因为物质将是像祂一样的天主——非受造、非生、无始无终。天主说：\Quote{我是元始！}但如果物质与祂同永恒，祂又怎能是元始？在同永恒者和同时代者之间没有先后顺序。那么，物质也是元始吗？\Quote{我，}主说，\Quote{独自铺张了诸天。}\BibleRef{依 44:24}但祂并非独自一人，因为那铺张穹苍的也同时是祂用以铺张的物质。当他主张物质是\Emph{永恒}的，且丝毫不侵犯天主的状况时，让他注意，我们不要以嘲笑回敬他：天主也同样是永恒的，且丝毫不侵犯物质的状况——二者的状况对祂们而言仍是共同的。因此，这一立场既未在物质方面受到攻击（它自身确实存在，只是与天主同在），也未在天主方面受到攻击（祂独自存在，却是与物质同在）。物质与天主同为元始，正如天主与物质同为元始；然而物质不能与天主比较，正如天主也不能与物质比较；物质与天主同为万物的创始者，并与天主同为万有之主权者。这样，\Emph{他提出天主}拥有物质的一部分，却又不是全部。因此，赫尔摩格内斯为天主保留了任何没有同样赋予物质的东西，以至于不是物质与天主被比较，而是天主与物质被比较。既然如此，既然那些我们宣称唯独属于天主的性质——永恒存在、无始无终、是元始、是独一、是万物的创始者——也同样适用于物质，我想知道物质拥有什么不同于天主且外在于天主的、因此而专属于它自己的性质，使它不能与天主相比较？那个拥有天主所有性质的存在，无须进一步比较就已被充分确定。\par

% structure: p0014 source=h2 target=section reason=relative-heading-rank; base=h2

\section{第七章 赫尔摩格内斯坚持他的理论，以便根据他自己的原则暴露其荒谬}

当他争论说物质小于天主、低于天主，因此与天主不同，并且因此不适合与那更伟大、更优越的天主比较时，我用这个原则反驳他：永恒而非受生的东西不可能有任何减少和低下，因为正是这一点使得天主本身如此伟大，低于且服从于无——不仅如此，比一切更伟大、更高超。因为，正如所有受生的、或终将终结的、因此不是永恒的东西，由于它们同时暴露于终结和开始，而具有与天主相悖的性质——我是指减少和低下，因为它们受生且受造——同样，天主也正是由于这个原因不受这些偶性影响，因为祂绝对非受生，也非受造。然而，物质的状况也是如此。因此，在两种永恒的存在（作为非受生非受造的存在）——天主和物质——中，由于它们共同状况的同一模式（两者都同样拥有那不容减少也不容服从的属性，即永恒性），我们断言两者都不比对方更小或更大，都不比对方更低或更高；相反，它们在伟大上相等，在崇高上相等，\Emph{并且}处于同一完整的、完美的福乐水平，永恒被认为是由此构成的。现在，我们不可在意见上像外邦人那样；因为他们被迫承认天主时，坚持在天主之下有其它的神。然而，神性没有等级，因为它是唯一的；如果它可以在物质中被发现——作为同样非受生、非受造和永恒——那么它必须同样存在于两者之中，因为它在任何情况下都不能低于自身。那么，赫尔摩格内斯将如何有勇气做出区分，从而将物质置于天主之下，将永恒置于永恒之下，将非受生置于非受生之下，将作者置于作者之下？因为物质敢于说：我也是元始；我也在万有之先；我从万有而出；我们平等，我们共存——两者都无始无终；两者都无作者、无天主。那么，那将我置于一个同时的、同永恒的能力之下的，是什么天主？如果那被称为天主的就是祂，那么我自己也有我自己的（神圣的）名字。要么我是天主，要么祂是物质，因为我们两者都是对方所不是的。因此，你难道认为他没有使物质与天主平等，尽管，确实，他假装物质低于天主？\par

% structure: p0016 source=h2 target=section reason=relative-heading-rank; base=h2

\section{第八章 根据他自己的原则，赫尔摩格内斯总体上使物质高于天主}

不仅如此，他甚至将\OriginalTerm{原始物质}{Matter}置于天主之上，反而使天主服从于它，因为他主张天主从原始物质中创造了万物。如果天主为了创造世界而从原始物质中汲取资源，那么原始物质便已被证明是更优越的，因为它为天主提供了成就其工程的手段；而天主也因此显然服从于原始物质，因为后者的本质对他是不可或缺的。没有人不需要他所使用的东西；没有人不服从于他所需要的东西，正是为了能够使用它。同样，没有人使用属于他人的东西而不低于其财产被使用的那一位；也没有人将自己的东西借给他人使用而不在这一点上高于借用其财产的那一位。依此原则，原始物质本身无疑并不需要天主，反而将自身借给需要它的天主——它丰富、充足且慷慨——借给一个我猜想是太渺小、太软弱、太无能的，无法从虚无中形成他所愿之物。它确实为天主提供了一项伟大的服务，即在此刻赐予他手段，使他得以被认识为天主，并被称为全能者——只是他不再是全能的了，因为他没有足够的能力从虚无中产生万物。诚然，原始物质也为自己谋取了一些东西——甚至使自己与天主一同被承认为天主的同等者，更是他的助手；只是有一个缺点，即只有\Person{赫尔莫杰内斯}发现了这个事实，除了那些哲学家——所有异端的始祖。因为先知们对此一无所知，宗徒们迄今也不知道，我想，连基督也不知道。\par

% structure: p0018 source=h2 target=section reason=relative-heading-rank; base=h2

\section{第九章　从\Person{赫尔莫杰内斯}的原则得出的若干不可避免却不可容忍的结论}

他不能说天主是作为原始物质的\OriginalTerm{主}{Lord}而使用它来进行创造工程的，因为对于与自己同等的实体，他不可能成为主。好吧，但或许这是一个源自他人意志的头衔，他所享有的是一种不稳定的持有，而非主权，而且\Emph{那}到了这样的地步，尽管原始物质是邪恶的，他仍然忍受使用一个邪恶的实体，这当然是由于他自身有限能力的限制，使他无力从虚无中创造，而非由于他的能力；因为如果作为天主，他实际上拥有对那他知道是邪恶的原始物质的主权，他本应先将其转化为善——作为它的\OriginalTerm{主}{Lord}和善的天主——以便他可以使用善的东西，而不是恶的。但既然他无疑是善的，却又不是主，他所凭借他所拥有的能力，表明了他不得不屈服于原始物质的条件，如果他是它的主，他本应改善它。这就是必须给\Person{赫尔莫杰内斯}的回答，当他主张天主是凭借其主权使用原始物质时——甚至是他对原始物质没有任何权利，当然是因为他没有亲自制造它。那么，按照你的说法，邪恶必然出自\Emph{天主}本身，因为他——我不说是邪恶的创造者，因为他并没有形成它，而是——邪恶的许可者，作为拥有统治权的那一位。如果确实原始物质根本不属于天主，因为它是邪恶的，那么，当他使用属于他人的东西时，要么是出于不稳定的头衔，因为他需要它，要么是出于暴力占有，因为他比它更强。因为获取他人财产有三种方式：通过权利、通过许可、通过暴力；换句话说，通过主权、通过源自他人意志的头衔、通过武力。既然主权是不可能的，\Person{赫尔莫杰内斯}必须选择哪一种适合于天主。那么，他是通过许可还是通过暴力，从原始物质中创造了万物？但事实上，天主难道不是更明智地决定什么都不创造，而不是仅仅凭借他人的容忍或通过暴力创造，而且还是用一种邪恶的实体？\par

% structure: p0020 source=h2 target=section reason=relative-heading-rank; base=h2

\section{第十章　\Person{赫尔莫杰内斯}荒谬地使天主陷入何等窘境；他无异于使天主成为邪恶的创造者}

即使\OriginalTerm{原质}{Matter}曾是完美的善，祂想到他人的财产（无论它多么善，借助它来达成祂的目的），难道不也同样有失体统吗？因此，祂为了自己的光荣，以这样一种方式创造世界，竟暴露出祂对属于另一个（而且并不善）的实体负有义务，这实在是够荒谬的。那么，\Person{赫谟根尼}问道，祂难道不该从无中创造万物，使邪恶之物本身也可归因于祂的意志吗？凭良心说，我们中间的异端者的盲目必定何等严重，以至他们那样争辩：要么他们坚持相信另一位至善的天主，因为他们认为造物主是邪恶的创造者；要么他们与造物主一同设立原质，为要从原质而非造物主引出邪恶。然而，根本没有哪个神祇能免于这种可疑的处境，以至于能避免看起来甚至像是邪恶的创造者，无论他是谁——我并非说确实制造了邪恶，而是容忍邪恶由某个作者、从某个来源制造出来。因此，应当立刻告诉\Person{赫谟根尼}（尽管我们将关于邪恶方式的区分推迟到另一处），他的这一计谋并未产生任何效果。因为请注意，天主如何被发现是邪恶的——若非创造者，至少也是纵容者：因为祂，以其极大的良善，在创造世界之前就已容忍原质中的邪恶，虽然作为良善者、邪恶之敌，祂本应纠正它。因为祂要么能纠正却不愿，要么愿意却是一位软弱的天主而不能。如果祂能却不愿，祂自己就是邪恶的，因祂偏袒了邪恶；如此，祂现在就使自己受到邪恶的指控，因为即使祂并未创造邪恶，但既然如果祂反对其存在它就不会存在，那么当祂拒绝意愿其不存在时，祂自己就必定已使其存在。还有什么比这更可耻的呢？当祂愿意那存在而祂自己又不愿创造时，祂实际上是违背了祂自己，因为祂既愿意那祂不愿制造的东西存在，又不愿制造那祂愿意存在的东西。仿佛祂所意愿的是善，而祂拒绝成为创造者的却是恶。祂因不创造而判定为恶的，又因允许其存在而宣告为善。祂容忍邪恶如善，而非根除它，从而证明自己是邪恶的推动者；若出于自己的意愿则是犯罪，若出于必要则是可耻。天主必须是邪恶的仆人或朋友，因为祂与邪恶在原质中交往——不，还从其中的邪恶成就了祂的工程。\par

% structure: p0022 source=h2 target=section reason=relative-heading-rank; base=h2

\section{第十一章  \Person{赫谟根尼}竭力将邪恶从天主转移到原质，却未能使其与自己的论证一致}

但是，说到底，\Person{赫谟根尼}用什么证据说服我们原质是邪恶的呢？因为他不可能不把那归咎于邪恶的东西称为邪恶。现在我们提出这一原则：永恒的实体不可能有减少和从属，以至于被认为比另一个同永恒的存在者低下。因此，我们现在断言，邪恶甚至不能与它相容，因为它不可能从属于任何事物，因它是永恒的。但是，由于在其他方面，显而易见，什么是永恒的，如天主是至高的善，由此祂单独是善——作为永恒，因此是善——作为天主，邪恶如何能内在于原质呢？原质（既然它是永恒的）必然被认为是至高的善。否则，如果那永恒的也被证明能容纳邪恶，那么这邪恶也能被加于天主而有损于祂；因此，他如此急于将邪恶从天主身上移开是没有充分理由的；因为邪恶必然与永恒的存在者相容，甚至通过与原质相容，正如\Emph{赫谟根尼所做的那样}。但是，按目前的论证，既然永恒的可以被视为邪恶，那么这邪恶必证明是不可战胜、不可克服的，因它是永恒的；在这种情况下，我们努力\BibleQuote{从你们中间除去那行恶的人} \BibleRef{格前 5:13}是徒劳的；在这种情况下，天主徒然给我们这样的命令和诫命；不仅如此，当天主确实要以不公正的方式施行惩罚时，祂设定任何审判都是徒然的。但另一方面，如果邪恶有终结的时候，当邪恶之首，魔鬼，\BibleQuote{进入那为他和他的天使所预备的永火} \BibleRef{玛 25:41}——先被\BibleQuote{投入深渊} \BibleRef{默 20:3}；同样，当\BibleQuote{天主子女的显扬} \BibleRef{罗 8:19}已经\BibleQuote{使受造物脱离} \BibleRef{罗 8:21}邪恶，这受造物曾\BibleQuote{屈服于虚妄} \BibleRef{罗 8:20}；当牲畜恢复其本性的天真与完整，与田野的野兽和平共处，当小孩子也与蛇嬉戏\BibleRef{依 11:6}；当圣父将祂的仇敌，即作恶者，置于祂圣子的脚下——如果这样，\Emph{终结}与邪恶相容，那么必然地，\Emph{开端}也与它相容；而原质也将证明有开端，因为它也有终结。因为凡归于邪恶账下的事物，都与邪恶的状况相容。\par

% structure: p0024 source=h2 target=section reason=relative-heading-rank; base=h2

\section{第十二章  争论方式的改变  接受\Person{赫谟根尼}的前提，以显示它们将他引入何种混乱}

来吧，让我们假定物质是恶的，不，甚至是极恶的，当然是出于其\Emph{本性}，正如我们相信天主是善的，甚至极善的，同样出于\Emph{本性}。现在，本性必须视为确定且不变的，就像在物质中固执地固定在恶中，在天主中则是不动不变地固定在善中。因为，显然，如果本性在物质中容许从恶变为善，那么在天主中它也可能从善变为恶。此时，有人会说：那么，\BibleQuote{难道不能从这些石头中给亚巴郎兴起子孙吗？}\BibleRef{玛 3:9} \BibleQuote{毒蛇的种类难道不能结出悔改的果实吗？}并且\BibleQuote{义怒之子}不能成为和平之子，如果本性是不变的吗？我的朋友，你提及这些例子是轻率的。因为那些因受生而存在的事物，如石头、毒蛇和人，与无始的物质并不相似；因为它们的本性既有起始，也就可能有终结。但请记住，物质已被决定为永恒的，作为非受造的、无始的，因此设想其本性是不变且不朽的；这正来自赫尔摩格内斯本人的观点，他以此反对我们，否认天主能够从自身创造任何东西，理由是永恒者不能改变，因为如果它不是永恒的，它就会失去——按此观点——它原本所是，并通过改变而成为它原本所不是。至于主，祂也是永恒的，（他主张）祂不能成为与祂永远所是者不同的任何东西。那么，我就采纳他这个确定的观点，并借此驳斥他。我以同样的责难谴责物质，因为出于它，尽管它是恶的——不，极恶的——善的事物被创造了，不，\BibleQuote{极善的}事物：\BibleQuote{天主看了，认为好，且祝福了它们}\BibleRef{创 1:21}——当然是因为它们极大的善；绝不是因为它们是恶的或极恶的。因此，物质中可以容许改变；既然如此，它就已经失去了永恒的状态；简言之，它的美丽在死亡中凋零。然而，永恒不能失去，因为它之所以是永恒，正是由于它免于失落。同样，它也不能改变，因为既然是永恒的，它就绝不可能改变。\par

% structure: p0026 source=h2 target=section reason=relative-heading-rank; base=h2

\section{第十三章 赫尔摩格内斯的另一个论据：物质中具有某种善。其荒谬性}

这里会出现一个问题：受造物如何从物质中产生出善来，而这些受造物是没有任何变化而形成的呢？在那是恶的、不，极恶的物质中，如何会有善的、不，极善的种子呢？诚然，\BibleQuote{好树不能结出坏果子}\BibleRef{玛 7:18}因为没有天主不是善的；坏树也不能结出好果子，因为没有物质不是极恶的。或者，如果我们同意在它\Emph{里面}有一些善的胚芽，那么它将不再有统一的本性（遍及它），即完全是恶的本性；反而（我们面对）一个双重本性，一部分善，一部分恶；问题又会出现：在一个同时是善和恶的主体中，黑暗与光明、苦与甜怎能和谐共存呢？再次，如果善恶这样截然不同的性质能够结合，并赋予物质一个产生这两种果实的双重本性，那么绝对的善事物就不再归功于天主，正如恶事物不归于祂，而是两种性质都属于物质，因为它们源于物质的属性。这样的话，我们对善事物不感谢天主，对恶事物也不怨恨，因为祂没有产生任何具有祂自己特性的事。由此将清楚证明，祂是服从于物质的。\par

% structure: p0028 source=h2 target=section reason=relative-heading-rank; base=h2

\section{第十四章 戴尔都良将对手逼入困境}

现在，如果有人也主张，尽管物质可能给祂提供了机会，但仍是祂自己的意志引导祂创造了善的受造物，因为祂在物质中发现了善——尽管这也是一个可耻的假设——然而，无论如何，当祂也同样从同一物质中产生恶时，祂就是物质的奴仆，因为当然，祂并非自愿产生这些，祂除了从一个恶的\Emph{根基}中完成创造外别无他法——无疑，作为善者，祂是不情愿的；也是出于必然，因为不情愿；并且是奴役的行为，因为出于必然。那么，哪个想法更值得推崇：祂出于必然创造了恶，还是出于自愿？因为如果出于物质，祂确实出于必然创造了它们；如果出于无，则出于自愿。因为你现在徒劳地努力避免使天主成为恶的作者；因为，既然祂用物质创造了万物，它们就必须归因于祂自己，即创造者，正是因为祂创造了它们。显然，祂从何处创造万物这个问题，等同于祂是否从无中创造万物；而祂从何处创造万物并不重要，只要祂从那里创造万物，并且那里给祂带来最大的光荣。现在，出于祂自己意志的创造比出于必然的创造给祂带来更多的光荣；换句话说，从无中创造比从物质中创造更为光荣。相信天主是自由的，即使是作为恶的作者，也比相信祂是奴隶更值得。权力，无论是什么，都比软弱更适合祂。如果我们甚至承认物质中没有任何善，而是主凭自己的大能产生了祂所产生的一切善，那么其他一些问题也会同样合理地出现。首先，既然物质中毫无善，\Emph{显然}善不是由物质受造的，其明确的理由就是物质并不拥有它。其次，如果\Emph{善}不是\Emph{受造}于物质，那么它必须是由天主受造的；如果不是由天主，那它必须是由无受造的。——因为根据赫尔摩格内斯自己的说法，这就是选择。\par

% structure: p0030 source=h2 target=section reason=relative-heading-rank; base=h2

\section{第十五章 真理：天主从无中创造了万物，被从对手的挣扎中解救出来}

倘若善不是从物质中产生的——因为恶性的物质中本无善——也不是从天主而来的（按\Person{赫尔莫革诺斯}的立场，没有任何东西能从天主产生），那么善就是从无中创造的，因为它既非来自物质也非来自天主。如果善是从无中形成的，为何恶不可以呢？不仅如此，如果任何物是从无中形成的，为何一切事物不能如此？除非说，天主的德能不足以产生\Emph{一切}事物，尽管它从无中产生了某物。或者，如果善是从恶的物质中发出的——因为它既非来自无也非来自天主——那么它必然来自物质的转变，这与\Emph{物质}作为\Emph{永恒存在}而被宣称的不可改变属性相违背。因此，关于善得以存在的来源，\Person{赫尔莫革诺斯}现在不得不否认其可能性。但善仍然必须从那些他否认其可能性的来源之一产生。如今，如果否认恶来自无，目的是为了否认恶是天主的作为（因为从祂的意志看，恶有太多理由被视为来自祂），而声称恶来自物质，以便恶成为那物的属性（假定恶是由那物的实体所造），那么，正如我所说，天主仍将被视为恶的作者；因为，本应藉其同样的权能和意志属性，从物质中产生一切善——或干脆产生善——然而祂却\Emph{不仅没有}产生一切善，甚至还产生了一些恶——当然，要么祂愿意恶存在（如果祂有能力使恶不存在），要么祂没有足够的力量使一切成为善（如果祂渴望如此却未能成就）；因为，无论主是出于软弱还是出于意志而成为恶的作者，都没有区别。否则，为什么在创造了善（仿佛祂本身就是善）之后，又产生了恶（仿佛祂在善上有所欠缺），因祂并未局限于只创造与祂自身一致的事物？在完成祂自己的作品之后，为什么还要劳神去处理物质，同时产生恶，以确保祂因善而被单独承认为善，并防止物质因（被造的）恶而被视为恶？如果恶没有吹拂，善本会更加兴旺。因为\Person{赫尔莫革诺斯}本人也驳斥了某些人的论点，这些人主张恶是必要的，以通过对比彰显善；这种理解必须来自对比。因此，这并非产生恶的理由；但如果必须为引入恶寻找其他理由，为何不能甚至从无中引入呢？因为正是这个理由会免除主被认为是恶的作者之责备，而如今这个理由却为恶的\Emph{存在}开脱（当主从物质中产生恶时）。如果存在这种开脱，那么问题就完全被逼到了一个角落，那些人不愿在此找到答案，他们不考察恶本身的理由，也不分辨应如何将恶归于天主或与天主分离，实际上却使天主遭受许多极其不值的诽谤。\par

% structure: p0032 source=h2 target=section reason=relative-heading-rank; base=h2

\section{第十六章 一系列两难困境：它们表明\Person{赫尔莫革诺斯}无法逃脱正统结论}

因此，就在这教义的门槛上——我可能要在别处详细讨论它——我明确地确立我的立场：善与恶两者都必须归于天主，祂用物质造了它们；或者归于物质本身，祂用物质造了它们；或者两者都归于两者共同，因为它们是相连的——既归于创造者，也归于祂所创造的那物质；或者（最后）一者归于一位，另一者归于另一位，因为在物质与天主之外没有第三位。如果两者都应归于\Emph{天主}，那么天主显然就是恶的作者；但天主既是善的，就不能是恶的作者。同样，如果两者都归于\Emph{物质}，那么物质显然就是善的母亲，但物质既是全然恶的，就不能是善的母亲。但如果两者都应当归于两者共同，那么在这种情况下，物质也将与天主相比较；两者将平等，在善与恶上处于同等地位。然而，物质不应与天主相比较，以免造成两位神。如果（最后）一者归于一位，另一者归于另一位——也就是说，善归于天主，恶归于物质——那么，一方面，恶不可归于天主，另一方面，善也不可归于物质。而且，天主通过用物质造出善与恶，就有了与物质一同造了\Emph{它们}。既然如此，我不知道赫尔谟根如何能逃避我的结论；因为他认为天主不能是恶的作者，无论祂如何从物质造出恶，无论是出于祂自己的意愿、出于必然，还是出于（事理的）理由。然而，如果祂是恶的作者，而祂是实际创造者，物质只是由于提供质料而与\Emph{祂}相关联，那么你就取消了\Emph{你}引入物质的理由。因为这并不减少其\Emph{真实性}，即通过物质，天主显示出自己是恶的作者，尽管\Emph{你}明确假定物质是为了防止天主显得是恶的作者。因此，既然原因被排除，物质也就被排除；那么毫无疑问，天主必定从无中造了万物。恶是否在其中，我们将在弄清楚什么是恶、以及你如今所认为的恶是否真的是恶时，再作讨论。因为对天主而言，祂出于自己的意愿、从无中造出这些，比出于另一个的预定（如果祂从物质造出它们，必是出于预定）更为相称。自由而非必然，适合天主的性格。我宁愿祂甚至自愿创造恶，也不愿祂缺乏阻止恶创造的能力。\par

% structure: p0034 source=h2 target=section reason=relative-heading-rank; base=h2

\section{第十七章　天主造化之工的真实性。你不能丝毫偏离它而不陷入荒谬}

这一规则是由独一天主的性质所要求的，祂之所以是独一的，除了作为唯一的天主外，别无他法；而祂之所以是唯一的，除了与祂没有其他东西共存外，别无他法。同样，祂将是首先的，因为万有在祂之后；万有在祂之后，因为万有出于祂；万有出于祂，因为它们出于无：所以理性与圣经一致，圣经说：\BibleQuote{谁曾知道主的心意？谁作过祂的谋士？祂与谁商议？谁曾向祂指示智慧与知识的道路？谁先给了祂，使祂得偿还呢？} 当然没有人！因为祂身边没有别的能力、物质或本性属于祂自己以外的。但如果祂是借助某部分物质进行创造，那么祂必定从那物质本身接受了对其秩序的设计和处理，作为\BibleQuote{智慧与知识的道路}。因为祂必须按照那物的性质和物质的本质来运作，而不是按照祂自己的意志；由于祂自己的意志，祂甚至必须依照不是祂自己、而是物质的本质来制造恶的事物。\par

% structure: p0036 source=h2 target=section reason=relative-heading-rank; base=h2

\section{第十八章　赞颂天主的智慧与圣言，天主藉此从无中造了万有}

若任何质料对天主创造世界是必要的，正如\Person{赫尔摩格内斯}所假设的，天主在自己的智慧中拥有一个更高贵、更适合的质料——这个质料不应以哲学家的著作来衡量，而应从先知的话语中习得。唯独这智慧知晓主的心意。因为\BibleQuote{谁能知道天主的事和天主内的事，除非是天主内的圣神？}\BibleRef{格前 2:11}如今，祂的智慧就是那位圣神。这圣神曾是祂的参谋，祂智慧与知识的道路。\BibleRef{依 40:14}天主藉着祂、同着祂创造了万物。\BibleQuote{当祂准备诸天时}，（圣经）如此说，\BibleQuote{我与祂同在；当祂坚固风云之上悠扬的云彩，当祂稳固天下源泉时，我与祂同在，与祂一同编联这一切。我是祂所喜悦的；每天我在祂面前欢乐：因为祂完成了世界就欢乐，并在人子中显示祂的喜悦。}\BibleRef{箴 8:27}如今，谁不更愿认同这智慧是万物的泉源和起源——事实上，是所有质料的质料，不受终结，状态不杂，运动不躁，形式不丑，而是自然、恰当、匀称、美丽，真是连天主也可能需要的，因为祂要求自己的，而非他人的？的确，一旦天主知道这智慧对创造世界是必要的，祂便立即创造祂，并在自身内生出祂。圣经说：\BibleQuote{主获得了我在祂工作之始，在祂造化的起头。在万世以前，祂就奠定了我；在祂造地以前，在群山安放之前，在丘陵之前，祂就生了我在深渊之先孕育了我。}因此，让\Person{赫尔摩格内斯}承认，天主的智慧本身被宣告为受生和受造的，特别原因是我们不应认为除天主外还有任何存在是无始无终、非受造的。因为如果那从起初就内在于主、从主而来、在主内的东西，尚且不是无始的——我指的是祂的智慧，它在天主的思想中开始为安排创造工程而运筹时，才受生和受造——那么，在主之外的东西，怎么可能没有起始呢？但如果这同一智慧是天主的圣言——作为智慧，并且（作为祂）万物都是藉祂造成的，同样，没有智慧就没有秩序——那么，除了圣父之外，有什么能比天主的独生子和首生者圣言更古老，并因此更高贵呢？更不用说无始者比受生者更强，非受造者比受造者更有力。因为不需要创造主赋予存在的，其等级远高于需要创造主使之存在的东西。按照这个原则，那么，如果恶果真是无始的，而天主子是受生的（因为，天主说：\Quote{我的心中吐出了我最卓绝的圣言}），我不太确定恶是否不会由善引入，强的由弱的引入，正如无始的由受生的引入。因此，基于此，\Person{赫尔摩格内斯}将质料置于天主之前，因为他将质料置于圣子之前。因为圣子就是圣言，并且\BibleQuote{圣言就是天主}\BibleRef{若 1:1}，以及\BibleQuote{我与父原是一体}\BibleRef{若 10:30}。但毕竟，也许圣子会足够耐心地容忍那（被\Person{赫尔摩格内斯}）置于祂之上、与圣父同等的东西！\par

% structure: p0038 source=h2 target=section reason=relative-heading-rank; base=h2

\section{第十九章 诉诸创造的历史。\Quote{起初}一词的真实含义，异端者诡异地曲解为荒谬的意义}

但我将诉诸梅瑟的原始文献，对方徒劳地试图借助它来支撑他们的猜想，当然，目的似乎是为了在这样一项探讨中显得得到那种不可或缺的权威的支持。他们找到机会，如异端者惯常所为，曲解某些词语的明显含义。例如，他们将\Emph{起初}一词，即天主创造天地之时，解释为仿佛意指某种有形的实体，被视为物质。然而，我们坚持每个词的正确含义，\Emph{并且说}\Emph{\OriginalTerm{开始}{principium}}意味着开始——这是一个适用于表示开始存在之物的术语。因为任何受造之物都不能没有开端，它的起始也不可能在它开始存在之前的任何其他时刻。因此，\Emph{\OriginalTerm{开始}{principium}}或开始，仅仅是一个起始的术语，而不是一个实体的名称。既然天地是天主的首要工程，并且由于祂首先创造了它们，以此在特殊方式下构成了祂万物的开端，先于一切，所以圣经有充分理由在其创造记录前加上这些话：\BibleQuote{在起初天主创造了天地}\BibleRef{创 1:1}；正如假若天主在其余一切之后才创造它们，就会说\Quote{最后天主创造了天地}一样。现在，如果开端是一个实体，那么终结也必定是有形的。无疑，一个实体可以是某种其他由它形成的事物的开端；例如，粘土是器皿的开端，种子是植物的开端。但是，当我们使用开端一词表示\Emph{起源}的意义，而非\Emph{秩序}的意义时，我们也不省略提及我们视为其他事物之起源的那个特定事物的名称。另一方面，如果我们这样说，例如：\Quote{起初，陶工制作了一个盆或一个水罐}，那么开端一词在此并不指示一种物质实体（因为我并未提及粘土，即\Emph{在此意义上}的开端，而只是指工作的\Emph{秩序}，意思是陶工先制作了盆和罐，在一切之前——然后打算制作其余的物品）。因此，开端一词所指涉的是工程的秩序，而非其物质的起源。我也可以用另一种方式解释\Emph{开端}一词，这并非不合适。希腊文中表示开始的词是\OriginalTerm{开始}{\GreekText{ἀρχή}}，它不仅包含秩序的优先性，也包含权力的意义；因此，王公和长官被称为\OriginalTerm{王公和长官}{\GreekText{ἀρχοντες}}。因此，在这种意义上，\Emph{开端}也可以被理解为王权的权威和力量。确实，天主是在祂超越的权威和力量中，创造了天地。\par

% structure: p0040 source=h2 target=section reason=relative-heading-rank; base=h2

\section{第二十章 短语\Quote{在起初}的含义——\Person{戴尔都良}将其与天主的智慧相连，并从中引出真理：创造并非出于先存的物质}

但在证明希腊词无非指\InnerQuote{起初}，而\Emph{起初}一词除了\Emph{初始}的意义外别无他意时，我们有那一位甚至承认这样一种起头，祂说：\Quote{主拥有了我，祂造化的开端，作为祂作为的开始。}\BibleRef{箴 8:22}因为万物都是藉天主的智慧造成的，由此可知，当天主\OriginalTerm{在起初}{in principio}——就是说，在起初——造成天地时，祂是在祂的智慧内造成它们。实际上，如果起初具有\Emph{物质}的意义，圣经就不会告诉我们天主\OriginalTerm{在起初}{in principio}创造了如此这般，而会更可能说\OriginalTerm{从起初}{ex principio}，即从物质中；因为祂不会\Emph{在}物质内创造，而是\Emph{从}物质中创造。然而，当提及智慧时，说\InnerQuote{在起初}是完全恰当的。因为祂首先在智慧内造了万物，因为祂在其中默想并安排祂的计划时，实际上已经完成了（创造之工）；即使祂有意从物质中创造，也仍会在祂的智慧中预先默想和安排时成就祂的创造，因为智慧事实上是祂作为的开端：这一默想和安排是智慧的首要运作，它藉默想和思想之行为为作为开辟道路。我甚至因此理由而主张圣经的权威：它既向我显示创造的天主和祂所造的工，却并未同样向我启示祂创造的来源。因为在每一作为中都有三要素：创造者、被造物以及造它的材料，因此，在关于作为的正确叙述中应提及三个名称——创造者的位格、所造之物的种类、以及形成它的材料。如果材料未被提及，而同时工作与工作的创造者都被提及，那么显然祂是从无中造出这工作。因为如果祂有任何可供操作的东西，它就会像（其他两项）一样被提及。最后，我将引用福音作为旧约的补充见证。现在，在其中更应当显示（如果有的话）天主造万物所用的材料，因为其中已经清楚启示了祂藉谁造万物。\Quote{在起初已有圣言}\BibleRef{若 1:1}——即，同样的起初，当然，天主在其中创造了天地\BibleRef{创 1:1}——\Quote{圣言与天主同在，圣言就是天主。万物是藉着祂造成的；凡受造的，没有一样不是藉着祂造成的。}\BibleRef{若 1:1}既然这里清楚告诉我们造物主是谁，即天主，以及祂造了什么，即万物，并藉谁造万物，即祂的圣言，那么叙述的顺序岂不应也声明万物由天主藉圣言而造所用的来源吗？如果万物真是从某种东西造的。因此，圣经无法提及不存在的东西；并且通过不提它，就清楚证明了没有这种东西：因为如果有，圣经就会提及。\par

% structure: p0042 source=h2 target=section reason=relative-heading-rank; base=h2

\section{第二十一章 对异端反驳的回应：圣经无须明确告诉我们世界由无中受造是多余的}

但是，你会对我说，如果你断定万物皆由无中受造，理由是并未告诉我们任何事物是由先在的物质所造，那么你要小心，免得被相反一方辩驳说，基于同样理由，万物皆由物质所造，因为也没有明确说任何事物是由无中受造。当然，有些论证很容易这样被反驳；但这并不意味着它们因此就可被接受，因为理由不同。因为我坚持认为，即使圣经没有明确宣告万物由无中受造——正如它避免（说它们）由物质（形成）——也没有必要特别指明万物由无中受造，正如如果有物质起源，就需要指明万物由物质受造。因为，对于由无中创造的事物，恰恰由于没有指明它由任何特定事物所造，就表明它由无中受造；并且，当根本没有指明它由什么所造时，就不存在它由某物所造的危险。然而，对于由某物所造的事物，除非明确宣告它由某物所造，否则就会有危险：首先，由于没有说它由什么所造，它看起来就像是由无中受造的；其次，如果它的本性使得它看起来肯定是由某物所造，那么同样有危险，即它似乎是由一种与恰当材料完全不同的材料所造，只要缺乏关于它由什么所造的声明。那么，如果天主不能从无中创造万物，圣经就不可能另外添加说祂从无中创造了万物：（这样的陈述就没有余地，）但它必定会以一切可能的方式告诉我们，祂从物质中创造了万物，因为物质必定是来源；因为一种情况即使没有明确陈述也完全可以理解，而另一种情况若不陈述就会存疑。\par

% structure: p0044 source=h2 target=section reason=relative-heading-rank; base=h2

\section{第二十二章 结论经圣经创造史的使用得以证实。\Person{赫尔莫根尼斯}面临增加圣经所宣告的灾祸}

圣神将圣经的这条规则确立到了如此程度：每当某物由某物造出时，祂总是同时提及所造之物和所由造之物。\BibleQuote{让大地，}祂说，\BibleQuote{生出青草、结种子的菜蔬和结果子的树木，各从其类，果子内含有种子，各从其类。就如此成了。大地就生出了青草、结种子的菜蔬，各从其类；并结果子的树木，果子内含有种子，各从其类。} \BibleRef{创 1:11} 又说：\BibleQuote{天主说：让水大量滋生有生命的动物，并让飞鸟在地上飞越穹苍天门。就如此成了。天主就创造了大鱼，以及水中滋生的一切有生命的动物，各从其类。} \BibleRef{创 1:20} 后来又说：\BibleQuote{天主说：让大地生出活的生物，各从其类：牲畜、爬虫和地上的野兽，各从其类。} 因此，如果天主在从已造之物中产生其他事物时，通过先知指明它们，并告诉我们祂是从何种来源产生了什么（尽管我们本可自己推测它们源自某个非虚无的来源，因为已有某些东西被造，它们似乎很容易由此被造）；如果圣神为教导我们付出了如此大的关怀，使我们知道万物由何产生，那么祂岂不也会同样妥善地告诉我们关于天和地的信息，指明祂是用什么造它们的，如果它们的原初是由某种物质构成的——祂似乎越是从无中造它们，事实上那时已造之物就越少，从那里祂似乎造了它们？因此，正如祂向我们指明祂从给定来源引出的事物之来源，同样，对于那些祂没有指明从何处产生的事物，祂（借由缄默）确认了我们的断言：它们是从无中产生的。\BibleQuote{在起初，天主创造了天地。} \BibleRef{创 1:1} 我崇敬祂圣经的丰富，在其中祂向我启示了创造者和受造物。此外，在福音中，我发现了创造者的仆人和见证者，即是祂的\OriginalTerm{圣言}{Word}。\BibleRef{若 1:3} 但是，万物是否都是由某种底层物质造成的，我至今在任何地方都未能找到。这样的陈述写在哪里，必须由\Person{赫尔摩格内斯}的店铺告诉我们。如果它任何地方都没有写，那么就让\Emph{它}畏惧那加给一切对\Emph{写成的话语}进行增删的人之\Emph{祸哉}吧。\BibleRef{默 22:18}\par

% structure: p0046 source=h2 target=section reason=relative-heading-rank; base=h2

\section{第二十三章 赫尔摩格内斯被追至另一段圣经。他的诠释之荒谬被揭露}

但他从下面的文字中提取了一个论证，那里写着：\BibleQuote{大地混沌空虚。} \BibleRef{创 1:2} 因为他将\Quote{大地}一词解为物质，因为由它造出的是大地。他对\Quote{是}一词也作了同样的解释，仿佛它指向那一直存在、非受生非受造之物。此外，它\Quote{混沌空虚}，因为他要物质曾经是无形、混乱的，并且没有经过工匠之手的修饰。现在我将逐一驳斥他的这些观点；但首先我想概括性地对他说：我们认为在这些词语中指涉的是物质。但是，圣经是否暗示，因为物质在万有之前就已存在，所以任何类似状态的东西是由它造出来的？完全没有。物质可能存在，如果它——或者更确切地说，如果赫尔摩格内斯——愿意的话。我可以说，物质可能存在，但天主可能并没有用它造出任何东西，要么因为祂不适合需要任何东西的帮助，要么至少因为祂没有被表明用了物质造了任何东西。因此，它的存在必是无因的，你会说。哦，不！当然不是无因的。因为即使世界不是由它造成的，但从那里孵出了一个异端，而且是一个特别无耻的异端，因为不是物质产生了异端，而是异端反而造出了物质本身。\par

% structure: p0048 source=h2 target=section reason=relative-heading-rank; base=h2

\section{第二十四章 大地并非如赫尔摩格内斯所愿意味着物质}

现在回到他以为物质被指涉的各个要点。首先我将考察那些术语。因为我们只读到其中一个，\Quote{大地}；另一个，即\Quote{物质}，我们并未遇到。我问，既然圣经中没有提及物质，\Quote{大地}这个术语怎能适用于它呢？这标志了另一种类的实体？如果\Quote{物质}又获得了\Quote{大地}的引申义，那么更有必要也提及\Quote{物质}，这样我才能确信\Quote{大地}和\Quote{物质}是同一个名称，从而不把该称号仅仅归于一种实体（作为其专有名称，并更由此为人所知）；或者，我无法（如果我有意的话）将其应用于物质的某个特定种类，而不是使之成为一切物质的通用术语。因为当一个事物没有专有名称而只有通用术语时，该术语所指的对象就越不明显，它就越能用于任何其他对象。因此，即使假设赫尔摩格内斯能向我们展示\Quote{物质}这个名称，他还必须向我们证明，同一对象还有\Quote{大地}这个别号，这样才能为它同时主张这两个称号。\par

% structure: p0050 source=h2 target=section reason=relative-heading-rank; base=h2

\section{第二十五章 关于创造史中提及两个大地的假设被驳斥}

因此他主张，在所述经文中向我们提出了两种地：一种是天主在起初创造的地；另一种是天主创造世界所用的物质，关于后者经上说：\Quote{大地还是混沌空虚。}\BibleRef{创 1:2}当然，若我问哪一种地最配称为\Quote{地}，我必被告知，受造的地是从那造它的物质得名，理由是：后裔从原初得名，比原初从后裔得名更合理。既然如此，另一个问题便呈现在我们面前：天主所造的这块地，是否应当从祂造它的物质得名？因为我从赫莫革尼斯和其他\OriginalTerm{物质主义异端者}{Materialist heretics}处得知：那一种地确实是\Quote{混沌空虚}，而我们的这块地却从天主那里同等获得了形状、美丽与匀称；因此受造的地与造它的物质是不同的。既然它已变成不同的事物，就不可能再与那物质共享其名，因为它已离开了那种状态。若\Quote{地}是（原始）物质的专有名称，我们这个世界的土地——它不是物质，因为它已成为别的事物——就不适合称为\Quote{地}，因为那名称属于另一物，与它的本性质不相干。但（你会告诉我）经历创造的物质，即我们的地，与它的原初既共享同类也共享名称。绝非如此。因为尽管陶瓮是由黏土制成，我却不再称它为黏土，而称它为陶瓮；同样，尽管\OriginalTerm{琥珀金}{electrum}由金和银合成，我却不会称它为金或银，而称它为琥珀金。任何事物一旦离开其本性，也就舍弃了它的名称——这种适宜性既为名称也为状态所要求。从那种作为物质的地到我们这块地，究竟发生了多么巨大的变化，从创世记中后者获得善的赞美就可见一斑：\Quote{天主看了认为好。}\BibleRef{创 1:31}而前者，按赫莫革尼斯所说，被视为一切邪恶的根源与原因。最后，若因为那一个是地，这一个就是地，那么为何那一个是物质，这一个却不是物质？其实，按此规则，天和所有受造物都应当被称为\Quote{地}和\Quote{物质}，因为它们都由物质组成。关于\Quote{地}这一名称，他用来指物质，我已说得够多。众所周知，这是元素之一的名称；因为首先由自然，然后由圣经教导我们如此，除非我们必须相信那个西勒诺斯，他在弥达斯王面前如此自信地谈论另一个世界——据特奥蓬波斯的记载。但同一作者告诉我们还有好几个神。\par

% structure: p0052 source=h2 target=section reason=relative-heading-rank; base=h2

\section{第二十六章 创造历史中遵循的方法，以回应赫莫革尼斯的曲解}

我们却只有一位天主，也只有一个地，就是天主在起初创造的地。\BibleRef{创 1:1}圣经在其开篇便打算叙述其秩序，首先告诉我们它被创造；接着说明这地是怎样的。同样，关于天，它首先告知其创造——\Quote{在起初天主创造了天。}\BibleRef{创 1:1}然后继续介绍其安排：天主如何将\Quote{穹苍以下的水与穹苍以上的水分开}，\BibleRef{创 1:7}并称穹苍为天——正是祂在起初所创造的那一个。同样，它（后来）论及人：\Quote{天主于是照自己的肖像造了人。}\BibleRef{创 1:27}接着揭示祂如何造他：\Quote{（上主）天主用地上的灰土形成了人，在他鼻孔内吹了一口生气，人就成了一个有灵的生物。}\BibleRef{创 2:7}这无疑是叙述的正确和适当方式：先有一个开场陈述，然后详述细节；先命名主题，再描述它。另一种看法多么荒谬：甚至在他还未预先提及他的主题——即物质——连名称都不给我们时，就突然公布了它的形状和状态，在提及它的存在之前就向我们描述了它的性质——指出了被造物的形状，\Emph{却}隐藏了它的名称！但我们的观点可信得多：圣经在适当描述了主题的形成并提及它的名称之后，才附加了它的安排。事实上，这些话语的意义多么充实完满：\Quote{在起初天主创造了天地；大地还是混沌空虚，}\BibleRef{创 1:1}——毫无疑问，正是天主所造的那同一个地，圣经当时刚刚论及它。因为那个\Quote{\Emph{但}}字被插入叙述中，像一个扣子，起着连词的作用，将两个句子不可分割地连接起来：\Quote{\Emph{但}大地。}这个词将思绪带回刚刚提及的那个地，并将意义约束在那里。除去这个\Quote{但}，联系就松开了；以至于\Quote{大地还是混沌空虚}这句话就似乎可能被理解为指任何别的地了。\par

% structure: p0054 source=h2 target=section reason=relative-heading-rank; base=h2

\section{第二十七章 其对手所沉溺的某些钻牛角尖的用词}

然而你接着称赞自己的眉毛，扬起头，用手指招呼，带着特有的轻蔑，说：那\Emph{\OriginalTerm{\Quote{是}}{was}}，似乎指向一种永恒的存在——当然，使其主词成为非受生、非受造的，并因此值得被设想为\OriginalTerm{物质}{Matter}。好吧，就我而言，我不会使用做作的抗议，只简单地回答：\Quote{\Emph{\OriginalTerm{是}{was}}}可以用于一切事物——甚至是一个被创造、被生、曾经不存在且不是你的\Quote{物质}的事物。因为对于一切存有的事物，无论其存在来自何处，无论是有开端还是无开端，\Quote{\Emph{\OriginalTerm{是}{was}}}这一词都将因其存在而得以使用。动词的第一时态适用于\Emph{定义}的，同一动词的后一时态在需要下降到\Emph{关系}时也适用。\Quote{\OriginalTerm{它是}{Est}}（它是）构成定义的本质部分，\Quote{\OriginalTerm{曾是}{Erat}}（曾是）构成关系的本质部分。这就是异端者的琐碎和诡辩，他们曲解并质疑最普通词语的简单含义。当然，一个重大问题在于：\Quote{那被造的地球\Emph{是}}！真正的讨论点在于：\Quote{\OriginalTerm{\Quote{无形}}{without form}和\OriginalTerm{\Quote{空虚}}{void}}这一状态更适合被创造的事物，还是更适合被创造所出的事物？这样，谓词（\OriginalTerm{是}{was}）可能属于主词（\OriginalTerm{那曾是的}{that which was}）所属的同一事物。\par

% structure: p0056 source=h2 target=section reason=relative-heading-rank; base=h2

\section{第二十八章 赫莫杰尼斯自相矛盾之处被揭露；创造史中的某些表述得到正确辩护}

但我们将证明，不仅这一状态适合我们的地球，而且它不适合赫莫杰尼斯所坚持的那另一种状态。因为，既然纯物质如此与天主共存，没有任何元素的介入（因为除了它自己和天主之外，尚不存在任何东西），它当然不可能是无形的。因为，尽管\Emph{\Person{赫莫杰尼斯}}主张黑暗是物质本性所固有的（这一观点我们将在适当的地方予以应对），但黑暗连人类都能看见（因为黑暗存在这一事实本身是明显的），对天主更是如此。如果它真是无形的，它的性质就根本无法被察觉。那么，赫莫杰尼斯是如何发现那物质是\Quote{无形}的，并且混乱无序，而作为无形之物，他的感官无法察觉呢？如果这一奥秘是由天主启示给他的，他应当向我们提供他的证据。我还想知道，那物质是否可以被描述为\Quote{空虚}。确实，那\Quote{空虚}的就是不完全的。同样确定的是，只有受造之物才可能不完全；当它尚未完全造成时是不完全的。当然，你承认这一点。因此，根本不曾受造的物质不可能是不完全的；而不完全的就不是\Quote{空虚}。它没有开端，因为它未被造，所以也不可能有任何空虚状态。因为空虚状态是开端的属性。相反，那受造的地球被恰当地称为\Quote{空虚}，因为它在受造之初、尚未完成之前，具有不完全的状态。\par

% structure: p0058 source=h2 target=section reason=relative-heading-rank; base=h2

\section{第二十九章 创造中宇宙秩序从混沌中逐步发展的美丽陈述}

天主确实按适当顺序完成了祂的一切工程；起初，祂仿佛先将它们以未成形的元素铺开，然后才将它们安排成完美的美。因为祂并没有一下子以太阳的光辉淹没光明，也没有一下子以月亮缓和的光线调和黑暗。祂并没有一下子以星座和星辰装饰天空，也没有一下子使海洋充满其繁多的怪物。祂也没有一下子赐给大地各种丰饶；而是首先赐予它存在，然后充满它，免得它被徒然创造。因为依撒意亚这样说：\BibleQuote{祂创造它并非徒然，祂塑造它以供居住。} \BibleRef{依 45:18} 因此，在它被造之后，等待其完美状态时，它是\Quote{混沌空虚}：实际上\Quote{空虚}，正是由于它没有形状（因为尚未完善到可见，同时也尚未具备其他性质）；而\Quote{没有形状}，是因为它仍然被水覆盖，仿佛被其孕育湿气的屏障所覆盖，我们的肉身由此产生，以一种与其相似的形式。因为达味正是此意：\Quote{大地属于上主，及其中的一切；世界和所有住在其间的：祂将大地奠基于海洋，建立于江河之上。} 正是当水被撤入其空洞深渊时，旱地变得显明，此前它一直为水所覆盖。然后它立即变得\Quote{可见}，天主说：\BibleQuote{让水聚集成一处，让旱地显现。} \BibleRef{创 1:9} \Quote{\Emph{显现}}，祂说，而不是\Quote{\Emph{被造}}。它早已被造，只是在不可见的状态下等待显现。\Quote{旱}，因为它将因与水分分离而成为这样，但仍然是\Quote{地}。\BibleQuote{天主称旱地为\Emph{大地}}，\BibleRef{创 1:10}，而非物质。因此，当它后来达到完美时，它不再被视为空虚，当天主宣告说：\Quote{让大地生出青草，结种子的菜蔬，各从其类，以及结果子的树木，果子内含有种子，各从其类。} 又说：\Quote{让大地生出活物，各从其类，牲畜、爬虫和地上的野兽，各从其类。} 如此，神圣的圣经完成了其完整的秩序。因为对于起初描述为\Quote{没有形状（看不见）和空虚}之物，它赋予了可见性和完成。此外，没有别的物质是\Quote{没有形状（看不见）和空虚}。因此，从此以后，物质必须是可见且完整的。所以我必须看见物质，因为它已成为可见。我也必须承认它是一个完成之物，以便能够从中收取结种子的菜蔬和结果子的树，并且由此产生的活物可以服务于我的需要。然而，物质无处可见，而大地却在此，呈现在我眼前。我见到它，我享受它，自从它不再是\Quote{没有形状（看不见）和空虚}以来。关于它，依撒意亚确实说过：\BibleQuote{上主如此说，祂创造了诸天，祂是塑造大地并造成它的天主。} \BibleRef{依 45:18} 祂所塑造的同一大地，也正是祂所造成的。那么，祂是如何塑造它的呢？当然是通过说：\BibleQuote{让旱地显现。} \BibleRef{创 1:9} 如果它先前不是不可见的，祂为什么命令它显现？\Emph{祂的目的}也是，通过使它可见从而适于使用，以免祂徒然造了它。因此，种种证据向我们表明，我们所居住的大地就是由天主创造并塑造的同一大地，而除了那被创造和塑造的以外，没有别的曾是\Quote{没有形状和空虚}。所以，\Quote{起初大地是混沌空虚}这句话适用于同一大地，即天主与诸天分开提及的那大地。\par

% structure: p0060 source=h2 target=section reason=relative-heading-rank; base=h2

\section{第三十章 另一段创世圣史，摆脱\Person{赫尔莫革涅斯}的曲解}

下列话语同样显然会证实\Person{赫尔莫革涅斯}的猜测：\BibleQuote{黑暗笼罩深渊，圣神在水面上运行；} \BibleRef{创 1:2} 仿佛这些混合的物质，为我们提供了他堆积如山的\Emph{物质}的论据。然而，如此分明的列举（正如我们在本文所见），分别指定了\Quote{黑暗}、\Quote{深渊}、\Quote{圣神}、\Quote{水}，这就排除了任何混淆或由此混淆导致的不确定意味。再者，当祂赋予它们各自的位置：\Quote{黑暗\Emph{在深渊的面上}}，\Quote{圣神\Emph{在水面上}}，祂否定了物质中的任何混淆；通过显示它们分立的位置，祂也显示了它们的区别。的确，最荒谬的是，物质被介绍给我们为\Quote{没有形状}，竟应通过如此多指示形状的词语来维持其\Quote{无形状}状态，却毫不暗示那混乱的物体是什么，它当然必须被认为是单一的，因为它是没有形状的。因为没有形状的东西是单一的；但即使是那没有形状的东西，当它由各种组成部分混合而成时，也必然有一个外观；但它没有任何外观，直到它从合并的许多部分中获得那唯一的外观。那么，物质要么本身具有这些特定的部分，即我们须从指示这些部分的词语来理解的——我是指\Quote{黑暗}、\Quote{深渊}、\Quote{圣神}、\Quote{水}——要么并不具有。如果它具有，它又怎能被介绍为\Quote{没有形状}？如果它不具有，它又如何被认识？\par

% structure: p0062 source=h2 target=section reason=relative-heading-rank; base=h2

\section{第三十一章　为圣经的创造叙述作进一步辩护，驳斥赫谟革内斯的虚妄观点}

然而，这一点也会被人钻空子，即圣经仅指着天和你的这个地，说天主在起初造了它们，但对上述具体部分却\Emph{未作任何说明}；因此，这些未被描述为受造的部分，就属于未成形的\OriginalTerm{物质}{Matter}。对此我们也必须予以答复。圣经若已宣告天地——作为创造中最崇高的工程——是由天主所造，它们当然拥有自己的特殊附属物，而这些附属物可被理解为包含在这最崇高的工程本身之中，那么圣经就已足够明确。当时在起初所造之天地的附属物，便是黑暗、深渊、圣神与水。因为深渊和黑暗是在地之下。既然深渊在地之下，黑暗在深渊之上，那么无疑，黑暗与深渊都在地之下。圣神与水也位于天之下。因为水覆盖着大地，圣神运行在水面之上，因此圣神与水都同样在大地之上。而凡在大地之上的，当然就在天之下。正如大地怀抱深渊与黑暗，同样，天也怀抱圣神与水，并将它们包容。其实，只提及作为整体的容器，而将其中所容之物视为部分包含在内，这并无新奇之处。假设我说这座城市建了一座剧场和一座竞技场，但舞台是某种样式，雕像在运河上，方尖碑矗立在其间；难道因为我未明确说明这些具体物件是由城市所造，它们就不是与竞技场和剧场一同被造的吗？我之所以不特意提及这些具体之物的形构，岂非因为它们在已说过被造之物中就已被暗示，并可被理解为内在于它们所容之物中？但此例或许因源自人的情况而显得牵强；我再举一例，它有圣经本身的权威。经上说：\BibleQuote{天主用地上的尘土造了人，在他鼻孔内吹入了生命的气息，人就成了一个有灵的生物。} \BibleRef{创 2:7}现在，虽然此处提到了鼻孔，但并未说明鼻孔是由天主造的；同样，在后面的经文中又提到皮肤、骨头、肉、眼睛、汗和血，却从未暗示它们是由天主创造的。赫谟革内斯将如何回答？人的肢体必须属于物质，因为它们未被特别提及为受造之物？还是它们包括在人的形造之内？同样，深渊、黑暗、圣神与水也\Emph{作为}天与地的肢体。因为在身体中，肢体被造；在身体中，肢体也被提及。没有元素不为其所容之物的肢体。但一切元素都包含在天与地之内。\par

% structure: p0064 source=h2 target=section reason=relative-heading-rank; base=h2

\section{第三十二章　创世纪中创造的记述是总括性的，却由旧约许多其他提及具体创造的段落所证实。进一步驳斥异议}

这是我应为维护我们面前的圣经而给出的回答，因为它在此似乎只陈述了天地之形成，仿佛它们是唯一被\Emph{造成}的物体。圣经不可能不知道，会有人从这些物体中立即理解它们的各个部分，因此它采用了这种简练的表达方式。但同时，它预见到会有愚蠢而狡猾的人，他们在真实含义上含糊其辞，并要求对各个部分也有一个描述其形成的词。正因为这样的人，圣经\Emph{才}在其他段落中教导我们各个部分的创造。你听见智慧说：\BibleQuote{但在深渊之前，我已受生，}\BibleRef{箴 8:24}为叫你们相信深渊也是\Quote{受生}的——即被造的——正如我们也生儿子，虽然我们\Quote{生他们}一样。深渊是被造还是被生，这无关紧要，只要它有一个开始；但若它附属于物质，就\Emph{不会}有开始。关于黑暗，主亲自通过依撒意亚说：\BibleQuote{我形成光，并创造黑暗。}\BibleRef{依 45:7}关于风，亚毛斯也说：\BibleQuote{祂加强雷霆，创造风，并向人宣示祂的基督；}\BibleRef{亚 4:13}由此表明，那风是被造的，它被视为与地的形成有关，在水面上运行，平衡、更新并赋予万物生气；而不是（如有些人所认为的）以\Emph{神}指天主自己，理由是\BibleQuote{天主是神，}\BibleRef{若 4:24}因为水不能承载它们的主；但祂是指构成风的那种神，正如祂通过依撒意亚所说：\Quote{因为我的神从我发出，我造了每一阵风。}同样，同一智慧论及水说：\BibleQuote{当祂使源泉坚固时，当祂使穹苍坚定时，我与祂一同塑造。}\BibleRef{箴 8:28}现在，当我们证明这些特殊事物是由天主创造的，尽管它们在创世纪中只是被提及，没有任何关于它们已被造成的暗示时，或许会从对方得到这样的答复：这些确实是被造的，但出自物质，因为梅瑟的陈述本身：\BibleQuote{黑暗笼罩在深渊的面上，天主的神在水面上运行，}\BibleRef{创 1:2}指的是物质，正如其他各处所有那些表明各个部分是从物质造成的圣经章节一样。那么，必然的结果是，正如地是由地构成的，同样深渊是由深渊构成的，黑暗是由黑暗构成的，风和海水是由风和海水构成的。并且，正如我们上面所说，物质不可能没有形式，因为它有特定的部分，这些部分是从它形成的——尽管作为分开的事物——除非它们并非分开，而是与它们所出自的那些事物完全相同。因为那些以相同名称陈述的特定事物，不可能是不同的；因为那样的话，天主的工作就可能显得无用，如果祂造了已经存在的事物；因为唯有当那些先前未被造成的事物出现时，才算是创造。因此，总结起来，要么梅瑟在写下\Emph{这些话}时指的是物质：\Quote{黑暗笼罩在深渊的面上，天主的神在水面上运行；}要么，既然这些创造\Emph{的部分}后来在其他经文中被表明是由天主造成的，那么它们就应当同样明确地被表明是出自物质，而按照你们的说法\Emph{，}梅瑟先前已提到过这物质；要么，\Emph{最后}，如果梅瑟指的是那些特殊部分，而不是物质，我想知道物质究竟\Emph{在何处}被指出过。\par

% structure: p0066 source=h2 target=section reason=relative-heading-rank; base=h2

\section{第三十三章 关于物质的真正教义的陈述。物质与天主创造世界的关系}

但是，虽然赫尔摩格尼斯在他自己的貌似合理的借口（因为他无力在天主的圣经中发现它）中找到了它，但对我们来说足够了：既确定万物都是由天主造成的，又丝毫不确定它们是由物质造成的。即使物质\Emph{先前}存在过，我们也必须相信它实在是天主所造的，因为我们坚持（同样）信德准则，即除了天主之外，没有什么是非受造的。到此为止，仍有争论的余地，直到物质受到圣经的检验，而不能证实其立场。整个结论如下：我发现没有一样东西是从无之外的任何东西造的；因为我发现被造的东西，我知道它\Emph{曾}不存在。凡是从某物造出的，其起源在于被造之物：例如，草、果实、牲畜和人的形体本身是从土地造出来的；同样，游泳和飞行的动物是从水产生的。产生这些受造物的原始质料，我可以说它们是它们的\Emph{质料}，但连这些质料也是天主创造的。\par

% structure: p0068 source=h2 target=section reason=relative-heading-rank; base=h2

\section{第三十四章 一个推定：万物由天主从无中创造，由万物最终归于无提供支持。证明这一归于无的圣经章节，从赫尔摩格尼斯指责其仅为比喻的说法中得以辩护}

此外，万物从无中受造的信念，将藉由天主那最终安排的眷顾铭刻在我们心中，这安排将促使万物归于无有。因为\Quote{诸天必要像卷轴一样被卷起}；不仅如此，它们还要与起初一同受造的大地一同归于无有。\BibleQuote{天地将要过去，} \BibleRef{玛 24:35} 主说。\BibleQuote{先前的天地已经过去了，} \BibleRef{默 21:1} \BibleQuote{再也找不到它们的地方，} \BibleRef{默 20:11} 因为当然，终结之物丧失了处所。同样，\Person{达味}说：\Quote{诸天，你手的化工，它们本身必要灭亡。你必将它们如同衣服一般更换，它们就被更换。}如今被更换，就是从它们原有的状态中坠落，在更换过程中失去那状态。\BibleQuote{星辰也要从天上坠落，如同无花果树被大风摇动而落下未熟的果子。} \BibleRef{默 6:13} \Quote{山岳在主面前要像蜡一般熔化}，就是说，\Quote{当祂起来剧烈震撼大地时。} \BibleRef{依 2:19} \BibleQuote{我要使水池干涸；} \BibleRef{依 42:15} 并且 \BibleQuote{他们必寻水，却寻不着。} \BibleRef{依 41:17} 连\Quote{海也不再有了}。如果有人竟以为这一切经文应当作属灵的解释，他仍不能剥夺它们那真实应验的事实，这些事件必照所写的成就。因为一切比喻性言辞必然源于真实事物，而非出于幻想；因为除非某物在比喻中实际是它所传递的那一事物，否则它不能传递任何属于自己的东西用作相似物。因此，我回到这个原则：凡从无中而来的，终必归于无有。因为天主不会用永恒的东西，即\OriginalTerm{物质}{Matter}，造出任何可朽之物；祂也不会从较优越之物造出较卑劣之物——较卑劣之物的本性更适合从较卑劣之物产生较优越之物，换言之，从可朽之物产生永恒之物。这正是祂甚至向我们的肉躯所应许的，并且祂愿意将祂自身德能与权能的这个保证放在我们内，好使我们相信，祂确实从无中唤醒了宇宙，仿佛它曾被浸没于死亡之中，这当然是指它先前不存在，以便进入存在。\par

% structure: p0070 source=h2 target=section reason=relative-heading-rank; base=h2

\section{第三十五章。\Person{赫尔谟根尼}关于物质及其属性的矛盾命题。}

关于涉及物质的其他所有点，虽然我们没有必要加以讨论（因为我们的首要之点是其存在的明显证据），但我们仍然必须进行讨论，就好像它确实存在一样，以便当这些关于物质的其他点彼此矛盾时，它的不存在就更加明显，同时也可使\Person{赫尔谟根尼}承认他自己的矛盾立场。他说，物质乍看之下似乎是\OriginalTerm{无形}{incorporeal}的；但当在\OriginalTerm{正确理性}{right reason}的\Emph{光照}下审视时，却发现它既非\OriginalTerm{有形}{corporeal}也非无形。你这正确理性到底是什么，它宣告的没有正确的，即没有确定的？因为，如果我没有弄错的话，一切事物必然要么有形要么无形（尽管我暂时可以承认甚至在实体事物中有某种无形性，尽管它们的实体本身就是特定事物的形体）；总之，在有形与无形之后，没有第三\Emph{种状态}。但如果有人主张，\Person{赫尔谟根尼}的这种正确理性发现了第三种状态，使物质既非有形也非无形，（我问）它在哪里？它是什么样的事物？它叫什么？它的描述是什么？它被理解为是什么？他的理性只宣告了这一点：物质既非有形也非无形。\par

% structure: p0072 source=h2 target=section reason=relative-heading-rank; base=h2

\section{第三十六章。以讽刺笔调揭露的关于物质及其偶性的其他荒谬理论。物质中的运动。\Person{赫尔谟根尼}关于它的狂妄臆想。}

但请看，他接着提出何等矛盾的说法（或者也许另有一个理由出现在他心中），当他宣称物质部分是有形体的，部分是无形的。那么，物质必须被视为（包含）两种状态，以便它可能既不具有其中之一？因为尽管那对立论宣称了这一点，它却会有形体和无形，而那对立论显然高于为自己的意见提供任何理由，正如那\Quote{另一个理由}也是如此。现在，他所说的物质中有形体的部分，是指身体由之受造的东西；而\Emph{物质}的无形部分，他是指其非受造的运动。他说：如果物质仅仅是一个身体，那么其中似乎就没有无形的东西，也就是（没有）运动；反之，如果它完全是无形的，就没有任何身体能由它形成。我们在这里拥有何等特殊的正确理由啊！\Person{赫尔莫根尼斯}，只要你画的草图与你推理的正确程度相当，那么没有哪个画家比你更愚蠢。因为谁会允许你把\Emph{运动}算作\Emph{物质}的一半呢？既然它不是实体性的事物，因为它不是有形体的，而是实体的一个依附属性（如果它甚至算是的话）？正如行动、推动、滑倒、跌倒一样，运动也是如此。当某物甚至自行运动时，其运动是推动的结果；但按照你的意思，当你把运动当作物质的无形部分时，它当然不是物质实体的任何部分。的确，一切事物都有运动——动物自身有，无生命之物由他物引起；然而我们不应因为人和石头既有身体又有运动，就说他们既有形体又无形体：\Emph{我们宁可说}，一切事物都具有单纯的形体性这一形式，那是实体的本质属性。如果任何无形的事件发生在它们身上，如行动、受动、功能、欲望，我们并不把这些算作事物的部分。那么，他如何设法将物质的一个\Emph{整体}部分归给\Emph{运动}呢？运动不属于实体，而属于实体的某种状态？这难道不是无可辩驳的吗？假设你突发奇想，要把物质描绘成不动的，那么静止在你看来会是其形态的一半吗？\Emph{当然不是}。同样，运动也不可能。不过，我将在别处有机会讨论运动。\par

% structure: p0074 source=h2 target=section reason=relative-heading-rank; base=h2

\section{第三十七章。关于物质的反讽式两难困境，以及妄加于它的各种道德属性}

我现在看到，你再次回到那个习惯向你宣告不确定之事的理由。正如你向我们介绍物质既非有形也非无形，你也声称它既非善也非恶；并且你在进一步以同样口吻论及它时说：\Quote{如果它是善的，既然它一直是善的，就不需要天主来安排它；如果它本性是恶的，它就不会容许变得更好，天主也永不会试图安排这样一种本性，因为祂的劳作将是徒然的。}这就是你的话，你若在别处也能记住这些话，避免自相矛盾就好了。然而，我们既已在一定程度上论及了与物质相关的这种善恶模糊性，我现在将只答复你摆在我们面前的这个命题和论证。我不再重复我的意见，即你本有责任明确说出物质是善还是恶，或是第三种状态；但（我必须指出）你在这里甚至没有坚持你先前选择的陈述。事实上，你收回了你的宣称——物质既非善也非恶；因为你暗示它是恶，当你说：\Quote{如果它是善的，它就不需要天主来安排；}同样，当你补充说：\Quote{如果它本性是恶的，它就不会容许任何变得更好的改变，}你似乎又暗示它是善的。这样你就赋予它一种与善恶的紧密关系，尽管你宣称它既非善也非恶。然而，为了驳斥你以为能巩固你的命题的论证，我在此主张：如果物质一直是善的，为什么它不会\Emph{仍然}想要变得更好呢？善的事物难道从不渴望、从不希望、从不觉得能够进步，从而将其善变为更善吗？同样，如果\Emph{物质}本性是恶的，为什么它不能被天主——作为更有力的存有——改变呢？正如祂能叫石头给亚巴郎兴起子孙来？\BibleRef{玛 3:9}毫无疑问，通过这些方式，你不仅将主与物质相比较，甚至将祂置于物质之下，因为你断言物质的本性不可能被祂控制和训练成为更好的。不过，尽管你在此不愿意承认物质本性是恶的，但在另一段中你将否认曾做出这样的承认。\par

% structure: p0076 source=h2 target=section reason=relative-heading-rank; base=h2

\section{第三十八章。\Person{赫尔莫根尼斯}关于物质及其某些附属物的其他推断，显为荒谬。例如，其所谓无限性}

关于原质之\Emph{位置}以及其\Emph{模式}的观察，均指向同一目标：即应对并驳斥你的乖谬主张。你将原质置于天主之下，因而自然为其指定了在天主之下的处所。故此原质是有位置的。既然有位置，它便处于位置之内；既处于位置之内，便为其所在之处所限定；既被限定，则有轮廓——你身为画师，深知轮廓乃一切可被勾画之物的边界。因此，原质不可能是无限的，因为它处于空间中，为空间所限定；既能为空间所确定，便可有轮廓。然而你却使之成为无限的，当你说：\Quote{它之所以无限，正因它始终存在。}倘若你的某位门徒想要声称你的意思是原质在时间上无限，而非在物质体量上无限，那么下文将表明你指的正是原质在体量上的无限性——它广大无边、不受限制。\Quote{因此，}你说，\Quote{它并非被整体塑造，而是在其部分\Emph{之中}塑造。}而当你既宣称原质在体量上无限，同时又为其指定位置、将其纳入空间与地方轮廓时，便自相矛盾。然而我仍不解：为何天主没有完全塑造它？除非祂要么无能，要么嫉妒。故我欲知那未被天主完全塑造的部分之\Emph{一半}，以便理解其整体为何。天主本应使之显明，作为远古的范本，以彰显其工程之荣耀。\par

% structure: p0078 source=h2 target=section reason=relative-heading-rank; base=h2

\section{第三十九章　后述之思辨被显明与赫谟根尼先前关于原质的首要原则相矛盾}

既然你认为更正确的做法是：让原质通过变化与位移来被限定；也让它能被理解，因为（如你所说）它被天主用作质料，基于其可转变、可变化、可分离。因为你说，它的变化显明它是不可分离的。在此你已偏离了你为天主的位格所规定的界限——你曾立下原则：天主不凭自己造它，因为祂不可能被分割，祂是永恒且永存的，因此是不变且不可分的。既然原质也被同一永恒性所衡量，无始无终，则它同样不可分割、不可变化，理由与天主相同。既然它在永恒性上与天主同享，则它必也在永恒的能力、法则与条件上分享。同样，当你说：\Quote{宇宙中所有事物同时皆拥有它的一部分，以便从部分知悉整体，}你自然是指那些由它产生、如今可见的部分。那么，宇宙中所有事物如何拥有（原质）——即从最初开始——既然如今可见之物在状态上与起初不同？\par

% structure: p0080 source=h2 target=section reason=relative-heading-rank; base=h2

\section{第四十章　无形原质作为天主美丽宇宙之起源实属不当。赫谟根尼推测仅用原质之一部分创造，并未改善其论证}

你说原质被改造得更好了——当然是从更坏的状态；\Emph{如此}你便使更好的成为更坏的翻版。万物曾混乱无序，如今井然有序；难道你还要说，从有序中产生无序？没有任何事物是另一事物的精确镜像；也就是说，并非其等同者。若有人以为在理发店的镜中看见自己像驴而非人，那除非他是认为无形无状的原质与如今在宇宙构造中被安排得美丽有序的原质相对应。如今世上有什么是无形的？当初在原质中又有什么是有形的，以致宇宙是原质的镜像？既然在希腊文中宇宙一词意为\OriginalTerm{装饰}{\GreekText{κόσμος}}，它如何能呈现未经装饰的原质之形，以至你能说整体由其部分可知？当然，那尚未成形的\Emph{部分}也属于整体；而你已经宣称\Emph{原质的全体}并未被用作创造的质料。因此，这部分粗糙、混乱、未整理的状态不能在精致、分明、有序的\Emph{被造物部分}中被识别出来——这些部分其实很难恰当地称为原质的部分，因为它们已脱离了原质的状态，在转变中与其分离了。\par

% structure: p0082 source=h2 target=section reason=relative-heading-rank; base=h2

\section{第四十一章　来自赫谟根尼的若干引文。其关于原质中运动以及善恶物质品质的思辨何其模糊不定}

我回到\Emph{运动}这一点上，好向你显示你在每一步上何等滑溜。物质中的运动是混乱的、杂乱的、动荡的。为此你把它比作一锅沸腾翻滚的热水。现在，在另一段中，你怎么会肯定另一种运动呢？因为当你想要表示物质既非善也非恶时，你说：\Quote{物质，作为（受造物）的基体，由于它具有均衡冲力的运动，既不十分趋向善，也不十分趋向恶。}如果它具有这种均衡冲力，它就不可能是动荡的，也不像锅中的沸水；它反倒会是平稳而有规律的，确实凭自身在善与恶之间摇摆，但并不倾向或趋向任何一边。它会像俗话所说，保持一种公正而精确的平衡。现在这并非不安；这并非动荡或不稳定；而是一种既不倾向任何一边的运动的规律性、平稳性和精确性。如果它一会儿向这边摇摆，一会儿向那边摇摆，并且更偏向某一特定边，那么它显然在该情况下就应受到不均匀、不平等和动荡的指责。此外，虽然\Emph{物质的}运动既不倾向于善也不倾向于恶，但它当然会在善与恶之间摇摆；因此从这一情况也可明显看出物质被限定在某种界限之内，因为它的运动既不倾向于善也不倾向于恶，由于它没有天生的偏向，就在两者之间摇摆，因此被限定在两者的界限之内。但事实上，你把善和恶都置于一个地域性的住所中，当你断言物质中的运动不倾向它们任何一个时。因为物质本是有地域性的，当它既不倾向这里也不倾向那里时，它就不倾向善和恶所在的地方。但当你把地点赋予善和恶时，你就使它们成为有形体的，因为你使它们成为有地点的；因为具有空间的东西必须首先有形体实体。事实上，无形体的事物不可能有自己独立的地点，除非当它们进入一个形体时，它们才能具有地点。但当物质不倾向于善和恶时，它是作为有形体的或有地点的\Emph{本质}而不倾向于它们。因此，当你坚持善和恶是实体时，你就错了。因为你给那些你指定地点的事物赋予了实体；但你指定地点，是因为你使物质中的运动保持平衡，距离两边相等。\par

% structure: p0084 source=h2 target=section reason=relative-heading-rank; base=h2

\section{第四十二章 进一步揭露\Person{赫尔莫琴斯}关于物质神圣属性的意见中的矛盾}

你随意而散漫地抛出了你所有的观点，免得它们因过于接近而显出自相矛盾。然而，我打算将它们收集起来加以比较。你声称物质中的运动是没有规律的，接着又说物质追求一种无形状的状态，然后在另一段中又说它渴望被天主整理。那么，那个追求无形的东西，想要被赋予形状吗？或者，那个想要被赋予形状的东西，却追求无形吗？你不愿意天主似乎与物质平等；然后你又说它和天主有共同的条件。\Quote{因为这是不可能的，}你说，\Quote{如果它跟天主没有任何共同之处，它就不能被祂整理。}但如果它与天主有任何共同之处，它就不需要被整理，因为它俨然是神性的一部分，由于条件相同；或者，连天主也有可能被物质整理，因为祂自己也与它有某种共同之处。现在你在此使天主受制于必然性，因为物质中有某种东西使天主给它赋予形状。你使两者都拥有这样一个共同属性：它们凭自身运动，并且永远在运动。你赋予物质的比赋予天主的少什么呢？在整个过程中，这种自由和永恒的运动都会表现出神性的共融。\par

只是在天主中运动是有规律的，在物质中运动是不规则的。但在两者中，神性的属性同样存在——两者都有自由和永恒的运动。同时，你赋予了物质更多的东西，它拥有这种自我运动的特权，而天主却没有被允许。\par

% structure: p0087 source=h2 target=section reason=relative-heading-rank; base=h2

\section{第四十三章 暴露并驳斥关于物质中的恶转变为善的其他矛盾}

关于运动，我还要说一点。遵循沸锅的\Emph{比喻}，你说物质中的运动在被调节之前是混乱的、不安宁的、因骚动过度而不可理解的。然后你接着说：\Quote{但它等待天主的调节，并因其不规则运动的迟缓而保持着不可理解的不规则运动。}刚才你赋予运动骚动，这里又赋予迟缓。现在观察一下你对物质的本性犯了多少错误。在前一段中你说：\Quote{如果物质本性是恶的，它就不会接受向善的改变；天主也绝不会对它施加任何整理的努力，因为祂的劳动会是徒然的。}因此你得出了两个意见：物质不是本性恶的，并且它的本性不能被天主改变；然后你忘记了它们，后来又得出这个结论：\Quote{但当它从天主那里接受调整并被归入秩序时，它放弃了它的本性。}既然它被转化为善，它当然是从恶转化而来的；如果通过天主的整理它放弃了恶的本性，那么它的本性就结束了；在调整之前它的本性是恶的，但在转化之后它可能已经放弃了那个本性。\par

% structure: p0089 source=h2 target=section reason=relative-heading-rank; base=h2

\section{第四十四章 关于天主与物质工作方式的奇特观点被揭露。异端者关于天主与物质空间关系的意见中的矛盾}

但我还应说明你如何使天主工作。你与哲学家们明显相左，但也与先知们不合。斯多葛派认为天主贯穿\OriginalTerm{质料}{Matter}，如同蜂蜜贯穿蜂巢。而你却断言，天主创造世界不是通过贯穿质料，而只是通过显现和接近它，如同美貌单凭显现就影响事物，磁石单凭接近就吸引事物。那么，天主形成世界与美貌刺伤灵魂或磁石吸引铁有何相似之处？因为即使天主向质料显现，祂并没有像美貌刺伤灵魂那样刺伤它；即使祂接近它，祂也没有像磁石与铁那样与之结合。然而，假定你的例子是合适的。那么，当然，天主是通过显现和接近质料而创造世界，祂是在显现和接近时创造了世界。因此，既然祂在此之前没有创造世界，祂就没有显现或接近过质料。那么，谁能相信天主没有向质料显现——而质料因其永恒性甚至具有同样的本性？或者祂曾远离它——即使我们相信祂无处不在、处处显现，万物都颂扬祂，甚至无生命和无形体之物，按（先知）达尼尔所言？\BibleRef{达 3:21} 多么浩大的地方，天主在创造世界之前如此远离质料，以致既未显现也未接近它！我想，祂是从遥远的地方旅行到它那里，一旦祂愿意显现和接近它。\par

% structure: p0091 source=h2 target=section reason=relative-heading-rank; base=h2

\section{第四十五章  结论。\Person{赫尔莫革尼斯}的言论与圣经关于创造的见证之间的对比。从无中创造，而非从\OriginalTerm{质料}{Matter}中创造}

但先知和宗徒们并非如此告诉我们，说世界是由天主仅仅通过显现和接近\OriginalTerm{质料}{Matter}而造成的。他们甚至没有提到任何质料，而是说智慧首先被设立，作为祂道路的开端，为祂的作为。\BibleRef{箴 8:22} 然后圣言被产生，\BibleQuote{万物是藉着祂造的；凡被造的，没有一样不是藉着祂造的。}\BibleRef{若 1:3} 的确，\BibleQuote{藉着上主的话，诸天造成；藉着祂口中的气，万象成列。}祂是上主的右手，\BibleRef{依 48:13} 其实是祂的双手，祂藉此工作并建造了\Emph{宇宙}。\BibleQuote{因为，}祂说，\BibleQuote{诸天是祢手的作为，}祂藉此\BibleQuote{量度诸天，以掌测地。}不要情愿用谄媚如此掩盖天主，以致争论祂仅凭显现和简单接近就产生了如此众多庞大的实体，而不是藉着祂自己的大能形成它们。因为耶肋米亚证明了这一点，他说：\BibleQuote{上主以自己的能力造了大地，以自己的智慧建立了世界，以自己的聪明展开了诸天。}\BibleRef{耶 51:15} 这些就是祂用以奋力创造宇宙的大能。如果祂劳苦，祂的荣耀更大。最后，第七日祂从祂的工程中安息了。两者都是按祂的方式。相反，如果祂只是通过显现和接近而创造了这个世界，那么祂在完成工程后，是否就不再显现和接近了？不，从世界被造之时起，天主反而更显著地显现，并处处可及。因此，你看，万物如何藉着那\BibleQuote{以能力造地、以智慧立世界、以聪明铺张诸天}的天主的作为而存在；不是仅仅显现或接近，而是运用祂心智、智慧、能力、聪明、圣言、圣神、大能的全能努力。如果祂仅仅通过显现和接近就成全了，这些东西对祂来说就不是必要的。然而，它们是祂的\BibleQuote{不可见的事}，按照宗徒所说，\BibleQuote{自从创造世界以来，藉着所造之物，就可以明明看见；}\BibleRef{罗 1:20} 它们\Emph{不是}一种难以名状的质料的部分，而是祂自身的可感证据。\BibleQuote{谁知道上主的心意？}\BibleRef{罗 11:34} 关于此，宗徒呼喊：\BibleQuote{深哉，天主的智慧和知识！祂的判断何其难测，祂的踪迹何其难寻！}这些话语表明什么更清楚的真理，岂不是万物都是从无中造出来的？它们除了天主以外，无人能寻见或探究。否则，如果它们在质料中可追溯或可发现，它们就是可探究的。因此，在\Emph{如此}程度上，质料并不先存已经显明（即使从这一点，即它不可能有被赋予的那种存在），在\Emph{那么}程度上，就证明了万物都是由天主从无中造出的。然而，必须承认，\Person{赫尔莫革尼斯}为质料描述了一种像他自己一样的状态——不规则、混乱、动荡，具有不确定、急躁和炽热的冲动——展示了他自己技艺的标本，并描绘了他自己的肖像。\par

\SourceRecord{Translated by Peter Holmes.；Ante-Nicene Fathers；Vol. 3.；Edited by Alexander Roberts, James Donaldson, and A. Cleveland Coxe.；Buffalo, NY: Christian Literature Publishing Co.,；1885.；<http://www.newadvent.org/fathers/0313.htm>.}
